\section{Externe Bibliotheken}

Für die Implementierung von TaskGrid+ wurden verschiedene Bibliotheken der
Python-Standardbibliothek sowie externe Frameworks eingesetzt. Tabelle
\ref{tab:bibliotheken} gibt einen Überblick über deren Zweck innerhalb des
Systems.

\begin{table}[H]
	\centering
	\begin{tabular}{p{3cm}p{13.75cm}}
		\toprule
		\textbf{Bibliothek} & \textbf{Verwendung} \\
		\midrule
		
		\texttt{grpc} &
		Implementierung der gRPC-Kommunikation zwischen Client, Dispatcher,
		Namensdienst und Workern. Bereitstellung von RPC-Schnittstellen sowie
		Erzeugung von Clients und Servern. \\
		
		\midrule
		
		\texttt{threading} &
		Nebenläufigkeit innerhalb des Systems. Verwendung für Locks,
		Heartbeat-Threads, Timeout-Überwachung und Hintergrundprozesse. \\
		
		\midrule
		
		concurrent.futures &
		Bereitstellung des \texttt{ThreadPoolExecutor} zur parallelen Verarbeitung
		mehrerer Tasks im Dispatcher. \\
		
		\midrule
		
		\texttt{queue} &
		Thread-sichere FIFO-Warteschlange zur Verwaltung ausstehender Tasks im
		Dispatcher. \\
		
		\midrule
		
		\texttt{http.server} &
		Implementierung des HTTP-Monitoring-Endpunkts
		(\texttt{GET /status}) des Dispatchers. \\
		
		\midrule
		
		\texttt{logging} &
		Strukturiertes Logging von Ereignissen, Zustandsübergängen und Fehlern
		zur Laufzeit. \\
		
		\midrule
		
		\texttt{json} &
		Serialisierung von Monitoring-Daten und Log-Ausgaben im JSON-Format. \\
		
		\midrule
		
		\texttt{dataclasses} &
		Definition kompakter Datenmodelle (z.\,B. Task-Objekte) mit automatisch
		generierten Konstruktoren und Hilfsmethoden. \\
		
		\midrule
		
		\texttt{enum} &
		Implementierung typisierter Aufzählungen, insbesondere für die
		Task-Zustände der Zustandsmaschine. \\
		
		\midrule
		
		\texttt{typing} &
		Bereitstellung von Typannotationen zur Verbesserung der Lesbarkeit,
		Wartbarkeit und statischen Codeanalyse. \\
		
		\midrule
		
		collections.defaultdict &
		Vereinfachte Verwaltung von Zählern und Zuordnungen, beispielsweise für
		die Round-Robin-Worker-Auswahl. \\
		
		\midrule
		
		\texttt{hashlib} &
		Berechnung kryptographischer Hashwerte (SHA-256) im Hash-Worker. \\
		
		\midrule
		
		\texttt{pathlib} &
		Plattformunabhängige Verarbeitung von Datei- und Verzeichnispfaden. \\
		
		\midrule
		
		\texttt{pytest} &
		Framework zur automatisierten Durchführung von Unit- und
		Integrationstests. \\
		
		\midrule
		
		\texttt{unittest.mock} &
		Erzeugung von Mock-Objekten zur Isolation einzelner Komponenten während
		der Tests. \\
		
		\midrule
		
		\texttt{tempfile} &
		Erzeugung temporärer Dateien und Verzeichnisse für Testzwecke. \\
		
		\midrule
		
		\texttt{io.StringIO} &
		Simulation von Datei- und Stream-Objekten im Speicher für Tests und
		Log-Validierungen. \\
		
		\bottomrule
	\end{tabular}
	\caption{Verwendete Bibliotheken und deren Funktion}
	\label{tab:bibliotheken}
\end{table}